% algtop_merged.tex
% A comprehensive introduction to algebraic topology, merging the
% detailed proofs from algtop_intro.tex with the full chapters on
% cohomology and further topics from algtop.tex. No information is lost.
\documentclass{article}
\usepackage[utf8]{inputenc}
\usepackage{amsmath, amssymb, amsthm}
\usepackage{tikz-cd}
\usepackage{hyperref}

\theoremstyle{definition}
\newtheorem{definition}{Definition}[section]
\newtheorem{example}{Example}[section]
\newtheorem{remark}{Remark}[section]
\newtheorem{theorem}{Theorem}[section]

\title{A Gentle Ascent into Algebraic Topology:\\From Holes to Homotopy}
\author{Compiled from lecture notes}
\date{}

\begin{document}
\maketitle
\tableofcontents

\section{What Algebraic Topology is About}
Algebraic topology tries to answer a simple question: \textbf{How can we tell whether two spaces are truly different?} 
A rubber sheet can be stretched and bent without tearing, so to a topologist a coffee mug and a doughnut are “the same” – they both have one hole.
We need \emph{invariants}: numbers, groups, or algebraic structures that stay unchanged under continuous deformations (homeomorphisms, or more generally homotopy equivalences).
The three fundamental invariants we will study are:

\begin{itemize}
    \item \textbf{Homology groups} $H_n(X)$ – detect “holes” of different dimensions.
    \item \textbf{Cohomology rings} $H^*(X)$ – a richer algebraic object that also remembers how holes interact.
    \item \textbf{Homotopy groups} $\pi_n(X)$ – classify maps from spheres into the space, extremely subtle.
\end{itemize}

We begin with the most geometric of these: homology.

\section{Homology: Counting Holes with Algebra}

\subsection{Chains and Boundaries}
To get hold of a topological space $X$, we approximate it by a simplicial complex – a union of points, edges, triangles, tetrahedra, etc.\ glued along faces.
An \textbf{oriented $p$-simplex} is the convex hull of $p+1$ ordered vertices $[v_0,\dots,v_p]$.
Orientation means that swapping two vertices changes the sign.

The \textbf{boundary} of a simplex is the alternating sum of its faces:
\[
\partial_p [v_0,\dots,v_p] = \sum_{i=0}^p (-1)^i [v_0,\dots,\widehat{v_i},\dots,v_p],
\]
where $\widehat{v_i}$ means the vertex is omitted.
We extend this operator linearly to formal sums of $p$-simplices, the \textbf{$p$-chains} $C_p(K)$.

\begin{theorem}[The boundary of a boundary is zero]
$\partial_{p-1} \circ \partial_p = 0$ for all $p$.
\end{theorem}
\textit{Proof idea.} Apply $\partial_{p-1}$ to the above sum: each face yields a sum of $(p-2)$-faces, and each such face appears twice with opposite signs because it can be obtained by omitting two vertices in either order.

\subsection{Homology Groups}
A $p$-chain $c$ with $\partial_p c = 0$ is called a \textbf{cycle} (it has no boundary).  
A $p$-chain of the form $\partial_{p+1} d$ is called a \textbf{boundary}.  
Because $\partial^2 = 0$, every boundary is a cycle. The $p$-th \textbf{homology group} is
\[
H_p(K) = \frac{\ker \partial_p}{\operatorname{im} \partial_{p+1}}.
\]
Intuitively, homology classes are cycles that are not themselves boundaries – that is, they detect “holes” that cannot be filled in.

\begin{example}[The circle $S^1$]
Approximated by a triangle, the chain $a+b+c$ in the picture is a cycle but not a boundary (there is no 2‑chain whose boundary is $a+b+c$). Hence $H_1(S^1) \cong \mathbb{Z}$.
\end{example}

\begin{example}[The torus $T = S^1 \times S^1$]
One 2‑dimensional hole (the interior) and two independent 1‑dimensional loops:
\[
H_0(T) = \mathbb{Z},\quad H_1(T) = \mathbb{Z}\oplus\mathbb{Z},\quad H_2(T) = \mathbb{Z}.
\]
\end{example}

\subsection{Relative Homology and the Long Exact Sequence}

Often we want to study a space modulo a subspace. For a pair $(K,L)$, we define \textbf{relative chains} $C_p(K,L) = C_p(K)/C_p(L)$. The boundary map descends to a map $\bar{\partial}_p$ on the quotient, and we obtain \textbf{relative homology} $H_p(K,L)$. Intuitively, chains in $L$ are ignored.

The relationships among $H_*(L)$, $H_*(K)$ and $H_*(K,L)$ are captured by a \textbf{long exact sequence}:
\[
\begin{tikzcd}
\cdots \arrow[r] & H_p(L) \arrow[r, "i_*"] & H_p(K) \arrow[r, "j_*"] & H_p(K,L) \arrow[r, "\partial_*"] & H_{p-1}(L) \arrow[r] & \cdots
\end{tikzcd}
\]
where $i:L\hookrightarrow K$ is inclusion, $j:K\hookrightarrow (K,L)$ is the natural map, and $\partial_*$ is a connecting homomorphism that “pushes a relative cycle onto its boundary in $L$”.

\subsubsection*{Construction and Proof of Exactness}
We want to understand why this sequence is exact, i.e. $\operatorname{Im} i_* = \ker j_*$, etc.
The proof is a standard application of the \textbf{zig‑zag lemma} applied to the short exact sequence of chain complexes
\[
\begin{tikzcd}
0 \arrow[r] & C_*(L) \arrow[r, "i"] & C_*(K) \arrow[r, "j"] & C_*(K,L) \arrow[r] & 0.
\end{tikzcd}
\]

\begin{enumerate}
    \item \textbf{Exactness of the chain sequence.}
    For each degree $p$, the map $i$ is injective and $j$ is surjective by definition, and $\operatorname{Im} i_p = \ker j_p$ because $C_p(K,L) = C_p(K)/C_p(L)$. Hence we have a short exact sequence of chain complexes.
    
    \item \textbf{The induced long exact sequence.}
    Applying homology to a short exact sequence of chain complexes always yields a long exact sequence of homology groups. The connecting homomorphism $\partial_* : H_p(K,L) \to H_{p-1}(L)$ is defined as follows:
    \begin{itemize}
        \item Take a relative cycle $\bar{c} \in C_p(K,L)$ (so $\bar{\partial} \bar{c}=0$). Lift $\bar{c}$ to an absolute chain $c \in C_p(K)$ via the surjection $j$. Then $\partial(c)$ maps to zero in $C_{p-1}(K,L)$ (because $\bar{\partial}\bar{c}=0$), so $\partial(c) \in \ker j_{p-1} = \operatorname{Im} i_{p-1}$. Hence there is a unique $b \in C_{p-1}(L)$ with $i_{p-1}(b)=\partial(c)$.
        \item One checks that $b$ is a cycle and its class $[b] \in H_{p-1}(L)$ is independent of the choice of lift. So we set $\partial_*[\bar{c}] = [b]$.
    \end{itemize}

    \item \textbf{Exactness at $H_p(K)$.}
    An absolute cycle $c$ satisfies $j_*[c]=0$ precisely when $j(c)$ is a relative boundary: $j(c) = \bar{\partial}(\bar{d})$. Lifting $\bar{d}$ to $d$, we get $j(c - \partial(d)) = 0$, so $c - \partial(d) \in \operatorname{Im} i$, which means $[c]$ comes from $H_p(L)$. Thus $\ker j_* \subseteq \operatorname{Im} i_*$. The opposite inclusion is trivial because $j\circ i =0$ on chains, so $j_*\circ i_* = 0$.
    
    \item \textbf{Exactness at $H_p(K,L)$ and $H_{p-1}(L)$.}
    Similar diagram chasing yields $\operatorname{Im} j_* = \ker \partial_*$ and $\operatorname{Im} \partial_* = \ker i_*$. The reader is encouraged to work through these zig‑zags explicitly.
\end{enumerate}

\begin{remark}
The exactness condition $\operatorname{Im} f = \ker g$ means that a homology class dies when it can be “filled” by a class from the previous group. This algebraic pattern repeats throughout topology.
\end{remark}

\subsection{Excision and the Mayer–Vietoris Sequence}
Two further powerful tools for computing homology are \textbf{excision} and \textbf{Mayer–Vietoris}.

\begin{theorem}[Excision]
If $U \subseteq A \subseteq X$ with $\overline{U} \subseteq \operatorname{int}(A)$, then the inclusion $(X\setminus U,\, A\setminus U) \hookrightarrow (X,A)$ induces isomorphisms $H_p(X\setminus U, A\setminus U) \cong H_p(X,A)$ for all $p$.
\end{theorem}
Intuitively, cutting out a subset that does not touch the “relative boundary” does not change relative homology.

\begin{theorem}[Mayer–Vietoris sequence]
Let $X = U \cup V$ be a union of two open sets (or, for simplicial complexes, two subcomplexes). Then there is a long exact sequence
\[
\begin{tikzcd}
\cdots \arrow[r] & H_p(U\cap V) \arrow[r, "\alpha"] & H_p(U)\oplus H_p(V) \arrow[r, "\beta"] & H_p(X) \arrow[r, "\Delta"] & H_{p-1}(U\cap V) \arrow[r] & \cdots
\end{tikzcd}
\]
\end{theorem}

\subsubsection*{Proof of Mayer–Vietoris (sketch)}
Consider the chain complexes of the two subspaces and their intersection. For each $p$, we have a short exact sequence
\[
\begin{tikzcd}
0 \arrow[r] & C_p(U\cap V) \arrow[r, "{x\mapsto (x,-x)}"] & C_p(U)\oplus C_p(V) \arrow[r, "{(a,b)\mapsto a+b}"] & C_p(U)+C_p(V) \arrow[r] & 0.
\end{tikzcd}
\]
The sum $C_p(U)+C_p(V)$ is a subcomplex of $C_p(X)$, and by a barycentric subdivision argument it computes the same homology as $C_p(X)$ (it is chain homotopy equivalent). Therefore we obtain a short exact sequence of chain complexes whose homology is the Mayer–Vietoris sequence, together with a connecting homomorphism $\Delta$ defined exactly as in the previous zig‑zag lemma. The map $\alpha$ is induced by the two inclusions, $\beta$ is the difference of the two inclusion‑induced maps, and $\Delta$ is the connecting map.

\subsection{Computing Homology of Spheres and Suspensions}
With these sequences in hand we can compute the homology of spheres.

\begin{theorem}
The homology of the $n$-sphere is
\[
H_p(S^n) \cong \begin{cases}
\mathbb{Z} & p = 0 \text{ or } p = n,\\
\{0\} & \text{otherwise}.
\end{cases}
\]
\end{theorem}

\textit{Proof by induction using Mayer–Vietoris.}
Decompose $S^n$ as the union of the upper hemisphere $U = S^n\setminus\{S\}$ and the lower hemisphere $V = S^n\setminus\{N\}$ (punctured at opposite poles). Each is homeomorphic to an open $n$-disk and thus contractible. Their intersection $U\cap V \simeq S^{n-1}$.
Now apply Mayer–Vietoris:
\[
\cdots \to H_p(U\cap V) \to H_p(U)\oplus H_p(V) \to H_p(S^n) \to H_{p-1}(U\cap V) \to \cdots
\]
For $p\ge 2$, the middle term vanishes (contractible spaces), giving isomorphisms
\[
H_p(S^n) \cong H_{p-1}(S^{n-1}).
\]
For $p=1$, we must check the connecting map. One finds $H_1(S^n)=0$ for $n>1$ and $H_0(S^n)\cong \mathbb{Z}$. Induction starting from $S^0$ (two points) yields the result.

\begin{example}[Suspension isomorphism]
If we apply Mayer–Vietoris to $\Sigma K$, the union of two cones $C_+K$ and $C_-K$ over $K$ intersecting exactly in $K$, we obtain isomorphisms
\[
H_{p+1}(\Sigma K) \cong H_p(K) \quad \text{for } p>0.
\]
This explains why the homology of a suspension shifts up by one dimension.
\end{example}

\subsection{Euler Characteristic}
If each chain group $C_p(K)$ has finite rank (say vector spaces over $\mathbb{R}$), the \textbf{Euler characteristic} is
\[
\chi(K) = \sum_{p=0}^{\dim K} (-1)^p \operatorname{rk} C_p(K).
\]
Thanks to the relation $\operatorname{rk} H_p = \operatorname{rk} Z_p - \operatorname{rk} B_p$ and a telescoping sum, one proves
\[
\chi(K) = \sum_{p=0}^{\dim K} (-1)^p \operatorname{rk} H_p(K),
\]
showing that $\chi$ is a topological invariant (even a homotopy invariant). For a sphere $S^n$, $\chi = 1+(-1)^n$; for a torus, $\chi = 0$.

\subsection{The Algebra of Homology: Tensor Products and Direct Sums}
Homology groups are abelian groups. The \textbf{direct sum} $G\oplus H$ simply pairs elements coordinate‑wise.  
The \textbf{tensor product} $G \otimes H$ is built from symbols $g\otimes h$ with the relation
\[
(g_1+g_2)\otimes h = g_1\otimes h + g_2\otimes h,\qquad g\otimes (h_1+h_2) = g\otimes h_1 + g\otimes h_2.
\]
These operations appear naturally when computing the homology of product spaces. Useful algebraic facts include:
\[
(G_1\oplus G_2)\otimes H \cong (G_1\otimes H)\oplus (G_2\otimes H),\qquad \mathbb{Z}\otimes G \cong G.
\]

\begin{theorem}[Künneth formula (simplified)]
If the coefficient group is a field, then
\[
H_p(X\times Y) \;\cong\; \bigoplus_{i+j=p} H_i(X) \otimes H_j(Y).
\]
For integer coefficients, a torsion correction term appears.
\end{theorem}

\section{Cohomology: The Dual Perspective}

\subsection{From Chains to Functions}
If homology is built from chains, cohomology is built from functions on chains.
Let $C_p(K)$ be the $p$-chains. Define \textbf{cochains} $C^p(K) = \operatorname{Hom}(C_p(K), \mathbb{Z})$ (or with coefficients in a ring $R$).
The \textbf{coboundary} operator $d_p : C^p(K) \to C^{p+1}(K)$ is the dual of $\partial_{p+1}$:
\[
(d_p \varphi)(c) = \varphi(\partial_{p+1} c).
\]
Because $d\circ d = 0$, we obtain \textbf{cohomology groups}
\[
H^p(K) = \frac{\ker d_p}{\operatorname{im} d_{p-1}}.
\]

\subsection{De Rham Cohomology: A Smooth View}
For smooth manifolds, the most natural avatar of cohomology comes from differential forms.
A \textbf{$p$-form} $\omega$ is an antisymmetric multilinear object that can be integrated over $p$-dimensional regions.
The \textbf{exterior derivative} $d\omega$ raises the degree by one and satisfies $d\circ d = 0$.

A form is \textbf{closed} if $d\omega = 0$, and \textbf{exact} if $\omega = d\eta$. The \textbf{de Rham cohomology} is
\[
H^p_{\text{dR}}(M) = \frac{\{\text{closed } p\text{-forms}\}}{\{\text{exact } p\text{-forms}\}}.
\]
Its key property is that it is isomorphic to singular cohomology with real coefficients, so it detects the same “holes”.

\begin{example}[Poincaré Lemma]
On a contractible manifold (like a star‑shaped domain in $\mathbb{R}^n$), every closed form is exact. Hence $H^p(M) \cong \{0\}$ for $p>0$.
\end{example}

\subsection{Wedge Product and the Cup Product}
The \textbf{wedge product} of two differential forms $\omega \in \Omega^p$, $\eta \in \Omega^q$ is a $(p+q)$-form defined by
\[
(\omega \wedge \eta)(v_1,\dots,v_{p+q}) = \frac{1}{p!\, q!}\sum_{\sigma} \operatorname{sgn}(\sigma)\, \omega(v_{\sigma(1)},\dots) \, \eta(v_{\sigma(p+1)},\dots).
\]
It satisfies $d(\omega\wedge\eta) = d\omega\wedge\eta + (-1)^p \omega\wedge d\eta$, so it descends to cohomology.
On the cohomological level this gives the \textbf{cup product}
\[
\cup : H^p(M) \times H^q(M) \longrightarrow H^{p+q}(M),\qquad [\omega]\cup[\eta] = [\omega\wedge\eta].
\]
This makes $H^*(M) = \bigoplus_{p\ge 0} H^p(M)$ into a graded ring. The ring structure is a powerful invariant: for example, it easily distinguishes $\mathbb{CP}^2$ from $S^2\vee S^4$ even though their homology groups are the same.

\subsection{Stokes' Theorem and Duality}
The crowning jewel of de Rham theory is \textbf{Stokes' theorem}:
\[
\int_\Omega d\omega = \int_{\partial\Omega} \omega.
\]
It encapsulates the fundamental theorem of calculus and all higher‑dimensional analogues. In the language of chains, it says that the evaluation pairing
\[
\langle [\omega], [c]\rangle = \int_c \omega
\]
between cohomology classes and homology classes is well‑defined and non‑degenerate over $\mathbb{R}$. This realizes $H^p(M;\mathbb{R})$ as the dual vector space of $H_p(M;\mathbb{R})$.

\begin{remark}[Cohomology vs.\ Homology]
Cohomology has advantages:
\begin{itemize}
    \item A ring structure (cup product) that captures subtle intersection information.
    \item Natural transformation laws under smooth maps (pullback).
    \item Links to global analysis, e.g.\ the interpretation of closed forms as solutions to differential equations.
\end{itemize}
\end{remark}

\subsection{Poincaré Duality via Morse Theory and Cellular Decompositions}
Poincaré duality is a deep symmetry between the homology and cohomology of a closed oriented manifold.  
For a compact, oriented \(n\)-manifold \(M\) without boundary, it asserts
\[
H^k(M;\mathbb{Z}) \;\cong\; H_{n-k}(M;\mathbb{Z}).
\]
Equivalently, the \(k\)-th cohomology group is isomorphic to the \((n-k)\)-th homology group.  
Morse theory provides a particularly geometric proof of this fact, by exhibiting a **cell decomposition**  and a **dual cell decomposition** that are the transpose of each other.

\begin{definition}[Morse function]
A smooth function \(f:M\to \mathbb{R}\) is called a \textbf{Morse function} if all its critical points are non‑degenerate (the Hessian matrix is invertible at each critical point).  
The \textbf{index} of a critical point is the number of negative eigenvalues of the Hessian, which equals the dimension of the **descending manifold** (the set of points that flow down to the critical point along \(-\nabla f\)).
\end{definition}

\textbf{Handle decomposition.}  
Around a critical point of index \(\lambda\), \(f\) looks like
\[
f(x_1,\dots,x_n) = f(p) - x_1^2 - \cdots - x_\lambda^2 + x_{\lambda+1}^2 + \cdots + x_n^2.
\]
The region below a regular value just above \(f(p)\) is obtained from the region below a regular value just below \(f(p)\) by attaching a \(\lambda\)-handle \(D^\lambda \times D^{n-\lambda}\).  
Thus, as we cross a critical point of index \(\lambda\), the manifold changes by attaching a **cell** of dimension \(\lambda\) (the core \(D^\lambda \times \{0\}\)).  
Passing through all critical values yields a cell decomposition of \(M\) where each critical point of index \(\lambda\) contributes one \(\lambda\)-cell.

\textbf{Morse–Smale complex.}  
If we choose a Riemannian metric generic with respect to \(f\), the descending manifolds of the gradient flow \(-\nabla f\) form a CW‑structure on \(M\) whose \(\lambda\)-cells are the descending manifolds of index‑\(\lambda\) critical points.  
The cellular chain complex \(C^{\text{Morse}}_*(M)\) then has a basis labeled by the critical points, and the boundary operator counts (with signs) the gradient flow lines connecting critical points of index difference \(1\).  
This chain complex computes the singular homology \(H_*(M)\).

\textbf{Dual decomposition.}  
If we replace \(f\) by \(-f\), the roles of ascending and descending manifolds swap.  
What was a \(\lambda\)-cell in the original decomposition becomes an \((n-\lambda)\)-cell in the new one, because the index of a critical point for \(-f\) is \(n - \text{index}_f\).  
The ascending manifold of \(p\) (for \(f\)) is precisely the descending manifold of \(p\) for \(-f\); it is a disk of dimension \(n-\lambda\) attached along its boundary to lower-dimensional cells.

Thus, the ascending manifolds form a **dual cell decomposition** of \(M\), with the property that the intersection number between a \(\lambda\)-cell and an \((n-\lambda)\)-cell of the dual decomposition is \(\delta_{ij}\) (if the cells come from the same critical point) or \(0\) (otherwise).  
Algebraically, this means that the chain complex built from the ascending manifolds is isomorphic to the dual complex \(\operatorname{Hom}(C^{\text{Morse}}_*, \mathbb{Z})\), shifted in degree by \(n\).  
Consequently,
\[
H^k(M) \;\cong\; H_{n-k}(M).
\]

\begin{remark}[Intuition]
The Morse function cuts the manifold into handles.  
The dual decomposition turns each \(\lambda\)-handle on its head, making it an \((n-\lambda)\)-handle in the reverse movie.  
Poincaré duality is the shadow of this geometric fact: the linear map that assigns to each cell its dual cell sets up an isomorphism between the cohomology of the original cellular cochain complex and the homology of the dual complex.
\end{remark}

This Morse‑theoretic proof not only establishes the duality but also equips it with a concrete geometric picture: homology classes are represented by descending cycles, cohomology classes by ascending cycles, and the pairing is given by their oriented intersection number.
\section{Homotopy Groups: Loops and Higher Spheres}
Homotopy groups detect higher-dimensional winding, not just holes. The exact sequences here are just as fundamental.

\subsection{The Fundamental Group}
\begin{definition}
For a pointed space $(X,x_0)$, $\pi_1(X,x_0)$ is the set of homotopy classes of loops based at $x_0$. The group multiplication is concatenation of loops.
\end{definition}

\begin{example}
\begin{itemize}
    \item $\pi_1(S^1) \cong \mathbb{Z}$: a loop that goes around the circle $n$ times gives the integer $n$.
    \item $\pi_1(S^n) \cong \{0\}$ for $n\ge 2$: any loop on the sphere can be shrunk to a point.
    \item $\pi_1(\mathbb{RP}^2) \cong \mathbb{Z}/2\mathbb{Z}$: the generator is a loop that cannot be shrunk, but going around twice is homotopically trivial.
    \item $\pi_1(T^2) \cong \mathbb{Z}\oplus\mathbb{Z}$: the two independent loops around the two holes.
\end{itemize}
\end{example}

\subsection{Higher Homotopy Groups and the Exact Sequence of a Pair}
For $n\ge 2$, $\pi_n(X,x_0)$ is defined using maps $f:(I^n,\partial I^n)\to (X,x_0)$. Equivalently, maps $S^n\to X$ preserving base points. The group operation is defined coordinatewise (and is abelian for $n\ge2$). There is a relative version $\pi_n(X,A,x_0)$ (or $\pi_n(X,A)$ for short) using maps from the $n$-cube that send one face to $A$ and the remainder of the boundary to $x_0$. These fit into the \textbf{long exact sequence of a pair}:
\[
\begin{tikzcd}
\cdots \arrow[r] & \pi_n(A) \arrow[r, "i_*"] & \pi_n(X) \arrow[r, "j_*"] & \pi_n(X,A) \arrow[r, "\partial"] & \pi_{n-1}(A) \arrow[r] & \cdots
\end{tikzcd}
\]

\subsubsection*{Construction and Exactness}
The maps are:
\begin{itemize}
    \item $i_*$ induced by inclusion $A\hookrightarrow X$,
    \item $j_*$ by regarding a map $(I^n,\partial I^n)\to (X,x_0)$ as a map $(I^n, I^{n-1}, J^{n-1})\to (X,A,x_0)$ where the face condition is satisfied,
    \item $\partial$ restricts a relative map to its face $I^{n-1}\times\{0\}$ that lands in $A$.
\end{itemize}

\textbf{Exactness at $\pi_n(X)$.} If $[f]\in\ker j_*$, then $f$ represents a relative homotopy class that is trivial, so there is a homotopy $H$ of maps $(I^n,\partial I^n)\to (X,x_0)$ relative to the face $I^{n-1}\times\{0\}$ (which sits in $A$). Actually, $j_*[f]=0$ means that $f$, viewed as a relative map, is homotopic relative to the distinguished face to the constant map at $x_0$. This homotopy provides a deformation of $f$ to a map whose image lies entirely in $A$, hence $[f]$ comes from $\pi_n(A)$. The detailed argument uses the fact that maps $I^n\to X$ are homotopic rel $\partial I^n$ to maps into $A$ if they can be deformed to lie in $A$ while fixing the boundary. Hence $\operatorname{Im} i_* = \ker j_*$.

\textbf{Exactness at $\pi_n(X,A)$.} The condition $\partial[\alpha]=0$ means that the restriction of $\alpha$ to the face $I^{n-1}$ is homotopic to the constant map inside $A$. One can extend this homotopy to a homotopy of the whole cube, obtaining a representative that is constant on that face, i.e.\ coming from an element of $\pi_n(X)$. Thus $\operatorname{Im} j_* = \ker \partial$.

\textbf{Exactness at $\pi_{n-1}(A)$.} Exercise.

\begin{remark}
The homotopy exact sequence of a pair does not require a simplicial structure; it works for any pair of spaces (with some niceness assumptions). It is fundamental for relating homotopy groups of a space and a subspace.
\end{remark}

\subsection{Fibrations and the Homotopy Exact Sequence}
A \textbf{fibration} $p: E \to B$ is a map satisfying the homotopy lifting property: any homotopy $H: X\times I \to B$ can be lifted to $\widetilde{H}: X\times I \to E$ if an initial lift is given. Every fibre bundle (locally trivial) is a fibration. For a fibration with fibre $F = p^{-1}(b_0)$, we have the \textbf{homotopy exact sequence}:
\[
\begin{tikzcd}
\cdots \arrow[r] & \pi_n(F) \arrow[r] & \pi_n(E) \arrow[r] & \pi_n(B) \arrow[r, "\partial"] & \pi_{n-1}(F) \arrow[r] & \cdots
\end{tikzcd}
\]
ending at $\pi_0(F)$, $\pi_0(E)$, $\pi_0(B)$ as sets.

\subsubsection*{Proof of Exactness (Sketch)}
The key ingredient is the homotopy lifting property. 
\begin{itemize}
    \item The map $\pi_n(F)\to \pi_n(E)$ is induced by inclusion.
    \item The map $\pi_n(E)\to \pi_n(B)$ is $p_*$. Its kernel consists of classes represented by maps $f: S^n\to E$ such that $p\circ f$ is nullhomotopic. Lifting the nullhomotopy gives a homotopy of $f$ into the fibre over $b_0$, so $[f]$ comes from $\pi_n(F)$. Thus $\operatorname{Im} i_* = \ker p_*$.
    \item The connecting map $\partial$: given $[g]\in\pi_n(B)$, represent it by $g: (I^n,\partial I^n)\to (B,b_0)$. The constant map at $e_0$ lifts to a map $\widetilde{g}$ defined on $I^{n-1}\times\{0\}$, and we extend it to a lift $\widetilde{G}$ of $g$ on the whole cube using the homotopy lifting property. The restriction of $\widetilde{G}$ to the remaining face $I^{n-1}\times\{1\}$ lands in the fibre $F$, giving an element of $\pi_{n-1}(F)$. This is independent of choices and defines $\partial[g]$.
    \item Exactness elsewhere follows from similar lifting arguments.
\end{itemize}

\begin{example}[Covering spaces]
The universal cover $\mathbb{R} \to S^1$, $t\mapsto e^{2\pi i t}$, is a fibration with fibre $\mathbb{Z}$ (discrete). Using the exact sequence and $\pi_n(\mathbb{R})=0$ for all $n$, we deduce $\pi_n(S^1)=0$ for $n\ge2$ and $\pi_1(S^1)=\mathbb{Z}$.
\end{example}

\subsection{The Hopf Fibration}
A remarkable map discovered by Hopf is the fibration
\[
S^1 \hookrightarrow S^3 \xrightarrow{\eta} S^2.
\]
Here $S^3 = \{(z_0,z_1)\in\mathbb{C}^2 \mid |z_0|^2+|z_1|^2=1\}$ and $\eta(z_0,z_1) = [z_0:z_1] \in \mathbb{CP}^1 \cong S^2$. The long exact sequence of this fibration gives
\[
\begin{tikzcd}
\cdots \arrow[r] & \pi_k(S^1) \arrow[r] & \pi_k(S^3) \arrow[r] & \pi_k(S^2) \arrow[r, "\partial"] & \pi_{k-1}(S^1) \arrow[r] & \cdots
\end{tikzcd}
\]
Since $\pi_k(S^1)=0$ for $k\ge2$, we obtain isomorphisms $\pi_k(S^3)\cong \pi_k(S^2)$ for $k\ge3$. In particular, for $k=3$, using $\pi_3(S^3)=\mathbb{Z}$ (by the Hurewicz theorem) we obtain
\[
\pi_3(S^2) \cong \mathbb{Z}.
\]
Thus there are non‑trivial maps $S^3 \to S^2$ that cannot be deformed to a constant map, the first example of a higher homotopy group being infinite even though no homology group would predict it.

\subsection{Projective Spaces and Orthogonal Groups}
Further identifications connect geometry and algebra:
\[
SO(3) \cong \mathbb{RP}^3 \cong S^3 / \{\pm 1\},\qquad
\mathbb{CP}^n \cong SU(n+1)/U(n).
\]
The Hopf map is then the quotient $S^3 \to S^3/S^1 \cong S^2$, where we view $S^3$ as the unit quaternions and $S^1$ as the complex numbers of unit norm inside them.

\subsection{The Hurewicz Theorem}
A deep relation between homotopy and homology is given by the \textbf{Hurewicz theorem}:
if $X$ is $(n-1)$-connected ($\pi_i(X)=0$ for $i<n$) with $n\ge 2$, then
\[
\pi_n(X) \cong H_n(X)
\]
is an isomorphism. For $n=1$, abelianization gives $H_1(X) \cong \pi_1(X)^{\text{ab}}$.
This theorem explains, for instance, why the first non‑vanishing homotopy and homology groups of a sphere coincide: $\pi_n(S^n) \cong H_n(S^n) \cong \mathbb{Z}$.

\section{Concluding Remarks}
We have covered the three pillars of basic algebraic topology:
\begin{itemize}
    \item \textbf{Homology}: cycles mod boundaries, counts holes; easy to compute with exact sequences.
    \item \textbf{Cohomology}: dual theory, has a ring structure; connected to differential forms and integration via de Rham cohomology.
    \item \textbf{Homotopy}: maps from spheres, much richer but harder to compute; related to homology by the Hurewicz theorem.
\end{itemize}

Along the way we encountered four families of exact sequences that are the workhorses of algebraic topology:
\begin{itemize}
    \item \textbf{Homology of a pair} and \textbf{Mayer–Vietoris} – for computing homology.
    \item \textbf{Homotopy of a pair} and \textbf{of a fibration} – for computing homotopy groups.
\end{itemize}
Their proofs, though elementary (diagram chasing or lifting arguments), demonstrate the power of categorical thinking. With them, one can compute invariants of a vast array of spaces.

The journey does not end here. Further topics include:
\begin{itemize}
    \item Higher homotopy theory (e.g.\ $\pi_4(S^3) \cong \mathbb{Z}/2\mathbb{Z}$).
    \item Spectral sequences for computing homology and homotopy.
    \item Sheaf cohomology and $K$-theory.
    \item Topological quantum field theories and the categorification of these invariants.
\end{itemize}

As with the study of adjunctions and Yoneda, the key is to see that \emph{every abstract invariant encodes a concrete geometric intuition}. The “holes” detected by homology, the “intersection numbers” recorded by the cup product, and the “winding” measured by homotopy are all facets of the same fundamental quest: to understand shape through algebra.

\end{document}